% SPDX-License-Identifier: Apache-2.0
% covers: NoReads
\datasheet{NoReads}{The read-data channel of a host whose client only
writes, held and never taken from.}

\dsfacts{\code{lib/parts/src/dma.rs}}{\code{txhdl\_parts::dma::NoReads}}%
{\code{//docs:datasheets}}

\subsection*{Function}

The other half of \code{NoBeats}. A store engine never reads, so its
host's read answers have nobody to go to, and a board joins every
channel it makes to a unit. \code{NoReads} holds that channel and takes
nothing from it.

\subsection*{Parameters}

\begin{center}\footnotesize
\begin{tabular}{@{}lp{0.7\textwidth}@{}}
\toprule
Parameter & Meaning \\
\midrule
\code{I} & The host's identifier width: the answer is \code{R<32, I>}.
It has no default, since the derive refuses one (issue 1044). \\
\bottomrule
\end{tabular}
\end{center}

The board has two: \code{NoReads<2>} for the Ethernet port's store
engine and \code{NoReads<1>} for the SD host's, whose shared host
carries one bit of identifier. The tables use \code{NoReads<2>}.

\dstables{NoReads}

\subsection*{Behaviour}

\begin{itemize}
\item It never takes an answer. In the netlist its ready is low for
ever.
\item Never taking is never a stall, because nothing is ever sent on
the channel: a client that issues no reads is answered with no read
data.
\item \code{idle} is its one register, and nothing sets it, as in
\code{NoBeats}.
\end{itemize}

\subsection*{Verification}

It has no example and no test of its own. It is in the board's
netlist, and the board tests in \code{cpu/vreteno/tests/board.rs} run
with it in place.

\subsection*{Used by}

\code{Board}, twice: as \code{snoreads}, \code{NoReads<2>}, on the
read-data channel of the Ethernet store engine's host \code{shost}, and
as \code{sdnoreads}, \code{NoReads<1>}, on the SD host's engines'
host \code{sdshost}.

\subsection*{Limits}

A host whose client does read must not be given one: the first answer
would stay in the channel, and the reads behind it would wait for
ever.
